\section{Introduction}
Open-source package dependencies distribute both functionality and trust. Earlier npm research established that dependency concentration and maintainer sharing create systemic points of failure~\cite{zimmermann2019}. Metadata weak-link signals also support proactive review~\cite{zahan2022}. The practical question is therefore not whether package graphs matter, but what a particular graph-derived queue measures and whether the artifacts used to construct it support that interpretation.

The repository studied here was developed around a proposal to forecast high-impact supply-chain entry points using dependency, maintainer, and public GitHub data. Its implemented collector stores npm package metadata, deps.dev dependency edges, workflow-directory observations, and imported advisories. A deterministic metric program combines maintainer portfolio size with the count of downstream packages associated with public workflows. A graph neural network is discussed in project documents but is not trained or evaluated by the available implementation.

This setting exposes a prerequisite often hidden by model-centric evaluation: a plausible score can be computed from inconsistent exports. We find that the compressed CSV and SQLite archives contain the same package population and identical maintainer and dependency metrics, while their workflow-derived metrics differ substantially. Consequently, using a shared population count or a common repository revision as a snapshot identifier is insufficient. A researcher can reproduce a reported CSV statistic faithfully while drawing conclusions inconsistent with the database in the same checkout.

We study three questions. \textbf{RQ1:} Are the supplied metric artifacts mutually consistent and reproducible from their stored graph and workflow observations? \textbf{RQ2:} How much do artifact choice, score weighting, and boundary ties change a fixed-budget defensive review queue? \textbf{RQ3:} What empirical claims remain supportable given the historical incident evidence and missing comparison population?

The resulting contribution has three parts. First, we provide a read-only audit that independently reconstructs all database metrics and compares every package across artifacts. Second, we quantify cross-artifact disagreement, queue overlap, and tie sensitivity without treating structural scores as compromise probabilities. Third, we specify a future benchmark whose time, population, outcome, and evidence requirements are explicit. Its graph model remains a hypothesis to test after those requirements are satisfied.

Our central result is a measurement result: \DBReachPositive{} of \PackageCount{} packages disagree on workflow reach, although every recomputed database metric agrees with its stored value. The two equal-weight score queues retain only \ScoreQueueIntersection{} of \ReviewBudget{} entries in common. A separate inventory of the historical companion yields \TemporalRecords{} event records across \TemporalFamilies{} family identifiers, \TemporalRetained{} retained source objects, and no eligible controls. These results support an artifact-consistency paper and a benchmark design, but cannot establish incident forecasting performance.
